\documentclass[a4paper]{article}
%\usepackage[utf8]{inputenc}
\usepackage[spanish, es-tabla, es-noshorthands]{babel}
\usepackage[table,xcdraw,dvipsnames]{xcolor}
\usepackage[a4paper, footnotesep=1.25cm, headheight=1.25cm, top=2.54cm, left=2.54cm,
 bottom=2.54cm, right=2.54cm]{geometry}
%\geometry{showframe}
 \usepackage[normalem]{ulem}
 \useunder{\uline}{\ul}{}

%VERIFICAR EL HEAD Y EL FOOT EN
%https://ctan.dcc.uchile.cl/macros/latex/contrib/geometry/geometry.pdf

%Paquetes varios:
\usepackage{verbatimbox}

%\usepackage{wrapfig}			%Wrap figure in text
\usepackage[export]{adjustbox}	%Move images
\usepackage{changepage}			%Move tables
\usepackage{todonotes}

\usepackage{tikz}
\usepackage{amsmath}
\usepackage{amsfonts}
\usepackage{amssymb}
\usepackage{float}
\usepackage[graphicx]{realboxes}
\usepackage{caption}
\usepackage{subcaption}
\usepackage{multicol}
\usepackage{multirow}
\setlength{\doublerulesep}{\arrayrulewidth}
%\usepackage{booktabs}

\usepackage{array}
\newcolumntype{C}[1]{>{\centering\let\newline\\\arraybackslash\hspace{0pt}}m{#1}}
%\usepackage[american]{circuitikz}
\usetikzlibrary{calc}
\usepackage{fancyhdr}
\usepackage{units} 

\usepackage{colortbl}
%\usepackage{sectsty}
%\usepackage{unicode-math}

%FONTS (IMPORTANTE): Compilar en XeLaTex o LuaLaTeX
\usepackage{anyfontsize}	%Font size
\usepackage{fontspec}		%Font type
%Si sigue sin andar comentar \usepackage[utf8]{inputenc}
%https://ctan.dcc.uchile.cl/macros/unicodetex/latex/fontspec/fontspec.pdf
%https://www.overleaf.com/learn/latex/XeLaTeX

%Path para imagenes para trabajar en subarchivos
\graphicspath{{../Resumen/}{../Referencias/}{../Apendice/}{../Descripción de la Empresa/}{../Tareas del Alumno/}{../Conclusiones/}{../Herramientas Empleadas/}}

%Definiciones de nuevos comandos y colores
%COLORES:
\definecolor{AzulFoot}{rgb}{0.682,0.809,0.926}	%RGB	%{174,206,235}
\definecolor{AzulInfo}{rgb}{0.180,0.455,0.710}	%RGB	%{46,116,181}
\definecolor{AzulTable}{rgb}{0.302,0.507,0.871}	%RGB	%{68,114,196}
\definecolor{PName}{rgb}{0.353,0.353,0.353}		%RGB	%{90,90,90}
\definecolor{mygreen}{rgb}{28,172,0} % color values Red, Green, Blue
\definecolor{mylilas}{rgb}{170,55,241}

%Change Font Size

% #1 = size, #2 = text
\newcommand{\setparagraphsize}[2]{{\fontsize{#1}{6}\selectfont#2 \par}}		%Cambia el size de todo el parrafo
\newcommand{\setlinesize}[2]{{\fontsize{#1}{6}\selectfont#2}}				%Cambia el font de una oración

%IMAGE IN TABLE:			%Ejemplo: \includeintable{.3}{ImagenesFactibilidad/pend}
\renewcommand\fbox{\fcolorbox{white}{white}}
\setlength{\fboxrule}{0pt}	%padding thickness
\setlength{\fboxsep}{4pt}	%border thickness
\newcommand{\includeintable}[2]{	
	\fbox{
		\begin{minipage}{#1\textwidth}
        	\includegraphics[width=\linewidth]{#2}
    	\end{minipage}
	}
}

%LINK IN REF
\newcommand{\reflink}[1]{		%LINK
	\href{#1}{#1}
}

%NOTAS:
\newcommand{\note}[1]{		%RED BIG NOTE (TODO)
	\begin{center}
		\huge{ \textcolor{red}{#1} }
	\end{center}
}

\newcommand{\lnote}[1]{{\fontsize{14}{6}\selectfont\textcolor{green}{#1}}}	%Notas pequeñas

\newcommand{\observacion}[2]{  \ifnumequal{1}{#1}{ { \todo[inline,backgroundcolor=red!25,bordercolor=red!100]{\textbf{Observación: #2}} } }{  }  }

\newcommand{\red}[1]{\textcolor{red}{#1}}

\newcommand{\TBD}{\textcolor{red}{(TBD) }}
\newcommand{\tbd}{\textcolor{red}{(TBD) }}

\newcommand{\TBC}{\textcolor{red}{(TBC) }}
\newcommand{\tbc}{\textcolor{red}{(TBC) }}

\newcommand{\quotes}[1]{``#1''}
\newcommand{\q}[1]{``#1''}

\newcommand{\ip}{192.168.0.10:1880}
\newcommand{\ipadmin}{192.168.0.10:1880/admin}

% Comandos para agregar elementos en tablas de acronimos y definiciones
\newcommand{\addacronym}[2]{\textbf{#1} & \begin{tabular}[l]{@{}l@{}}#2\end{tabular} \\ \hline}

% tabItem
\newcommand{\tabitem}{~~\llap{\textbullet}~~}


\usepackage{hyperref}
\hypersetup{
    colorlinks=true,
    linkcolor=black,
    filecolor=magenta,      
    urlcolor=AzulInfo,
    citecolor=AzulInfo,    
}

%Configuración del header y del footer:
\usepackage{etoolbox}
\pagestyle{fancy}
\fancyhf{}
\rfoot{\thepage}
\renewcommand{\footrulewidth}{4pt}
\renewcommand{\headrulewidth}{0pt}
\patchcmd{\footrule}{\hrule}{\color{AzulFoot}\hrule}{}{}

%Código en el informe
%% IMPORTANTE:
% Verificar que esté \usepackage[dvipsnames]{xcolor}

%\usepackage{listingsutf8}
\usepackage{listings}

\renewcommand{\lstlistingname}{Código}

%LSTSET: Pone un recuadro y contador de linea en el codigo
\newcommand{\boxstyle}{
	\lstset{
		basicstyle=\sffamily\color{black},
		frame=single,
		numbers=left,
		numbersep=5pt,
		numberstyle=\color{gray},
		showspaces=false,
		showstringspaces=false
	}
}

\newcommand{\defaultstyle}{
	\lstset{
		basicstyle=\sffamily\color{white},
		frame=none,
		numbers=none,
		showspaces=true,
		showstringspaces=true
	}
}

\lstdefinelanguage{Kotlin}{
  captionpos=b,
  comment=[l]{//},
  commentstyle={\color{gray}\ttfamily},
  emph={filter, first, firstOrNull, forEach, lazy, map, mapNotNull, println},
  emphstyle={\color{OrangeRed}},
  identifierstyle=\color{black},
  keywords={!in, !is, abstract, actual, annotation, as, as?, break, by, catch, class, companion, const, constructor, continue, crossinline, data, delegate, do, dynamic, else, enum, expect, external, false, field, file, final, finally, for, fun, get, if, import, in, infix, init, inline, inner, interface, internal, is, lateinit, noinline, null, object, open, operator, out, override, package, param, private, property, protected, public, receiveris, reified, return, return@, sealed, set, setparam, super, suspend, tailrec, this, throw, true, try, typealias, typeof, val, var, vararg, when, where, while},
  keywordstyle={\color{NavyBlue}\bfseries},
  morecomment=[s]{/*}{*/},
  morestring=[b]",
  morestring=[s]{"""*}{*"""},
  ndkeywords={@Deprecated, @JvmField, @JvmName, @JvmOverloads, @JvmStatic, @JvmSynthetic, Array, Byte, Double, Float, Int, Integer, Iterable, Long, Runnable, Short, String, Any, Unit, Nothing},
  ndkeywordstyle={\color{BurntOrange}\bfseries},
  sensitive=true,
  stringstyle={\color{ForestGreen}\ttfamily},
}

\lstdefinelanguage{Swift}
{
  morekeywords={
    open,catch,@escaping,nil,throws,func,if,then,else,for,in,while,do,switch,case,default,where,break,continue,fallthrough,return,
    typealias,struct,class,enum,protocol,var,func,let,get,set,willSet,didSet,inout,init,deinit,extension,
    subscript,prefix,operator,infix,postfix,precedence,associativity,left,right,none,convenience,dynamic,
    final,lazy,mutating,nonmutating,optional,override,required,static,unowned,safe,weak,internal,
    private,public,is,as,self,unsafe,dynamicType,true,false,nil,Type,Protocol,
  },
  morecomment=[l]{//}, % l is for line comment
  morecomment=[s]{/*}{*/}, % s is for start and end delimiter
  morestring=[b]", % defines that strings are enclosed in double quotes
  breaklines=true,
  escapeinside={\%*}{*)},
  numbers=left,
  captionpos=b,
  breakatwhitespace=true,
  basicstyle=\linespread{1.0}\ttfamily, % https://tex.stackexchange.com/a/102728/129441
}

\definecolor{keyword}{HTML}{BA2CA3}
\definecolor{string}{HTML}{D12F1B}
\definecolor{comment}{HTML}{008400}

\newcommand{\swiftstyle}{
	\lstset{
  		language=Swift,
  		inputencoding=utf8x,
		extendedchars=\true,
	  	basicstyle=\ttfamily,
	  	showstringspaces=false, % lets spaces in strings appear as real spaces
  		columns=fixed,
  		keepspaces=true,
  		keywordstyle=\color{keyword},
  		stringstyle=\color{string},
  		commentstyle=\color{comment}
	}
}


%Como usarlo:

%\begin{lstlisting}[caption={Simple code listing.}, label={lst:example1}, language=Kotlin]
%// this is a simple code listing:
%println("hello kotlin from latex")
%\end{lstlisting}

%Si se corta en 2 páginas distintas:

%\vspace{1mm}
%\noindent{\begin{minipage}{\linewidth}
%\begin{lstlisting}[...]
%...
%\end{lstlisting}
%\end{minipage}}




\usepackage{titlesec}		%Para hacer las subsubsubsections

%Colores a los nombres de las secciones:
%\sectionfont{\color{AzulInfo}}  % sets color of sections
%\subsectionfont{\color{AzulInfo}}
%\subsubsectionfont{\color{AzulInfo}}

%PICTURES AND TABLE INDEX:
\newcommand{\Section}[1]{ \section{#1} 
	\phantomsection \setcounter{figure}{0} \setcounter{table}{0} \setcounter{lstlisting}{0}
		\renewcommand{\thetable}{\arabic{section}.\arabic{table}}
		\renewcommand{\thefigure}{\arabic{section}.\arabic{figure}}
		\renewcommand{\thelstlisting}{\arabic{section}.\arabic{lstlisting}}
}

\newcommand{\Subsection}[1]{ \subsection{#1}
	\phantomsection \setcounter{figure}{0} \setcounter{table}{0} \setcounter{lstlisting}{0}
		\renewcommand{\thetable}{\arabic{section}.\arabic{subsection}.\arabic{table}}
		\renewcommand{\thefigure}{\arabic{section}.\arabic{subsection}.\arabic{figure}}
		\renewcommand{\thelstlisting}{\arabic{section}.\arabic{subsection}.\arabic{lstlisting}}
}

\newcommand{\Subsubsection}[1]{ \subsubsection{#1} 
	\phantomsection \setcounter{figure}{0} \setcounter{table}{0}  \setcounter{lstlisting}{0}
		\renewcommand{\thetable}{\arabic{section}.\arabic{subsection}.\arabic{subsubsection}.\arabic{table}}
		\renewcommand{\thefigure}{\arabic{section}.\arabic{subsection}.\arabic{subsubsection}.\arabic{figure}}
		\renewcommand{\thelstlisting}{\arabic{section}.\arabic{subsection}.\arabic{subsubsection}.\arabic{lstlisting}}
}

%Definición de subsubsubsection:
\titleclass{\subsubsubsection}{straight}[\subsection]

\newcounter{subsubsubsection}[subsubsection]
\renewcommand\thesubsubsubsection{\thesubsubsection.\arabic{subsubsubsection}}

\titleformat{\subsubsubsection}
  {\normalfont\normalsize\bfseries\color{AzulInfo}}{\thesubsubsubsection}{1em}{}	%Color de subsubsubsection
\titlespacing*{\subsubsubsection}
{0pt}{3.25ex plus 1ex minus .2ex}{1.5ex plus .2ex}

\makeatletter
\renewcommand\paragraph{\@startsection{paragraph}{5}{\z@}%
  {3.25ex \@plus1ex \@minus.2ex}%
  {-1em}%
  {\normalfont\normalsize\bfseries}}
\renewcommand\subparagraph{\@startsection{subparagraph}{6}{\parindent}%
  {3.25ex \@plus1ex \@minus .2ex}%
  {-1em}%
  {\normalfont\normalsize\bfseries}}
\def\toclevel@subsubsubsection{4}
\def\toclevel@paragraph{5}
\def\toclevel@paragraph{6}
\def\l@subsubsubsection{\@dottedtocline{4}{7em}{4em}}
\def\l@paragraph{\@dottedtocline{5}{10em}{5em}}
\def\l@subparagraph{\@dottedtocline{6}{14em}{6em}}
\makeatother

\setcounter{secnumdepth}{4}
\setcounter{tocdepth}{4}

%Subsubsubsection:
\newcommand{\Subsubsubsection}[1]{ \subsubsubsection{#1} 
	\phantomsection \setcounter{figure}{0} \setcounter{table}{0} \renewcommand{\thetable}{\arabic{section}.\arabic{subsection}.\arabic{subsubsection}.\arabic{subsubsubsection}.\arabic{table}} \renewcommand{\thefigure}{\arabic{section}.\arabic{subsection}.\arabic{subsubsection}.\arabic{subsubsubsection}.\arabic{figure}}
}

%Tamaño, color e identación de sección, subsección, subsubsección y subsubsubsección:
%La identación de las subsecciones está tambien en Index-cfg.tex para el toc, lot y lot en el index
\titleformat{\section}[block]{\fontsize{16}{6}\selectfont\bfseries\color{AzulInfo}}{\thesection.}{1em}{} 
\titleformat{\subsection}[block]{\hspace{2.5em}\fontsize{13}{6}\selectfont\color{AzulInfo}}{\thesubsection}{1em}{}
\titleformat{\subsubsection}[block]{\hspace{3.5em}\fontsize{12}{6}\selectfont\color{AzulInfo}}{\thesubsubsection}{1em}{}
\titleformat{\subsubsubsection}[block]{\hspace{4em}\fontsize{11}{6}\selectfont\color{AzulInfo}}{\thesubsubsubsection}{1em}{}

%Pone las refrencias en el indice
\usepackage[numbib, nottoc, notlot, notlof]{tocbibind}

%Pone toc, lof y lot en colores y elijo el titulo de estos
\addto\captionsspanish{
	\renewcommand\contentsname{Contenidos}
	\renewcommand\listfigurename{Lista de Figuras}
	\renewcommand\listtablename{Lista de Tablas}
}

%Agrega TOC al indice
\renewcommand{\tableofcontents}{
	\stepcounter{section}
	\addcontentsline{toc}{section}{\protect\numberline{\thesection}\textbf{Contenidos}}
	\tableofcontents
}

%Agrega LOF al indice
\renewcommand{\listoffigures}{
	\stepcounter{section}
	\addcontentsline{toc}{section}{\protect\numberline{\thesection}\textbf{Lista de Figuras}}
	\listoffigures
}

%Agrega LOT al indice
\renewcommand{\listoftables}{
	\stepcounter{section}
	\addcontentsline{toc}{section}{\protect\numberline{\thesection}\textbf{Lista de Tablas}}
	\listoftables
}

%Indices: cambio la separación de los numeros para que entren tablas y fotos
\usepackage{tocloft}
\setlength{\cftfignumwidth}{1.35cm}  % change numwidth from figures in lof
\setlength{\cfttabnumwidth}{1.35cm}  % change numwidth from tables in lot
\renewcommand{\cfttoctitlefont}{\Large\bfseries\color{AzulInfo}}
\renewcommand{\cftloftitlefont}{\Large\bfseries\color{AzulInfo}}
\renewcommand{\cftlottitlefont}{\Large\bfseries\color{AzulInfo}}

%Coloca lineas punteadas a las seciones en el TOC
\renewcommand{\cftsecleader}{\cftdotfill{\cftdotsep}}

%Items con bullets y no cuadrados
\renewcommand{\labelitemi}{\textbullet }


\begin{document}

\def\PrimerObs{1}
\def\SegundaObs{0}

%%%%%%%%%%%%%%%%%%%%%%%%%
%		Caratula		%
%%%%%%%%%%%%%%%%%%%%%%%%%

\setmainfont{Avenir LT Std 55 Roman}
\begin{titlepage}

\begin{tikzpicture}[remember picture, overlay, black, line width = 0.5pt]
	\coordinate (a) at (-2cm,2cm);
	\coordinate (b) at (17cm,-25.5cm);
	
	\coordinate (ap) at (-2.1cm,2.1cm);
	\coordinate (bp) at (17.1cm,-25.6cm);
	
	\draw[] (a) -| (b);
	\draw[] (a) |- (b);
	
	\draw[] (ap) -| (bp);
	\draw[] (ap) |- (bp);
	
	%footnotesep=1.25cm, headheight=1.25cm, top=2.54cm, left=2.54cm, bottom=2.54cm, right=2.54cm

\end{tikzpicture}

\begin{figure}[H]
	\includegraphics[width=0.3\linewidth, right]{./Utils/ITBA_1}
\end{figure}

\vspace*{1.5cm}

\noindent \textbf{\setlinesize{11}{INSTITUTO TECNOLÓGICO DE BUENOS AIRES - ITBA}}

\noindent \textbf{\setlinesize{11}{ESCUELA DE INGENIERÍA Y TECNOLOGÍA}}

\vspace*{4cm}

\center{
\setlinesize{24}{ \textbf{SISTEMA INTEGRAL DE MONITOREO DE FAUNA SILVESTRE} }

\vspace*{1.5cm}

%\setlinesize{24}{ \textbf{Subtítulo del trabajo (cuando corresponda)} }
}

\vspace*{1.5cm}

\begin{figure}[H]
\begin{adjustwidth}{-1cm}{}
\begin{tabular}{llr} 
	\textbf{AUTORES:} & \textbf{Mechoulam, Alan}  &  \textbf{(Leg. N}$\mathbf{^o}$ \textbf{58438)}\\
	 & \textbf{Lambertucci, Guido Enrique} & \textbf{(Leg. N}$\mathbf{^o}$ \textbf{58009)} \\
	 & \textbf{Rodriguez Turco, Martín Sebastian} & \textbf{(Leg. N}$\mathbf{^o}$ \textbf{56629)} \\
	 & \textbf{Londero Bonaparte, Tomás Guillermo} & \textbf{(Leg. N}$\mathbf{^o}$ \textbf{58150)} \\
 &  & \\
 &  & \\
	\textbf{DOCENTES}: & \textbf{Orchessi, Walter} & \\
	 & \textbf{Pingitore, Ricardo} & \\
	 & \textbf{Ugarte, Alejandro} & \\
\end{tabular}
\end{adjustwidth}
\end{figure}

\vspace*{0.5cm}
{\noindent \textbf{TRABAJO FINAL PRESENTADO PARA LA OBTENCIÓN DEL TÍTULO DE INGENIERO ELECTRÓNICO}}
\vspace*{1.5cm}

\center{\textbf{BUENOS AIRES}}

\end{titlepage}


\begin{titlepage}

\begin{tikzpicture}[remember picture, overlay, black, line width = 0.5pt]
	\coordinate (a) at (-2cm,2cm);
	\coordinate (b) at (17cm,-25.5cm);
	
	\coordinate (ap) at (-2.1cm,2.1cm);
	\coordinate (bp) at (17.1cm,-25.6cm);
	
	\draw[] (a) -| (b);
	\draw[] (a) |- (b);
	
	\draw[] (ap) -| (bp);
	\draw[] (ap) |- (bp);
	
	%footnotesep=1.25cm, headheight=1.25cm, top=2.54cm, left=2.54cm, bottom=2.54cm, right=2.54cm

\end{tikzpicture}

%\vspace*{4cm}
%\noindent \textbf{TRABAJO FINAL PRESENTADO PARA LA OBTENCIÓN DEL TÍTULO DE INGENIERO ELECTRÓNICO}
%
%\vspace*{5cm}
%
%\center
%
%\textbf{BUENOS AIRES}

\vspace*{1cm}

\centering{\textbf{PRIMER / SEGUNDO CUATRIMESTRE, 202X}}


\end{titlepage}


\setmainfont{Calibri}
\begin{titlepage}
\begin{figure}[H]
	\centering
	\includegraphics[width=0.5\linewidth]{./Utils/ITBA_2}
\end{figure}

\vspace*{1.5cm}

\center
{ \setlinesize{48}{Proyecto final de Ingeniería Electrónica} }

\vspace*{1cm}

{ \textcolor{PName}{\setlinesize{22}{Sistema Integral de Monitoreo De Fauna Silvestre}} }

\vspace*{3cm}

\begin{tabular}{llr} 	
\textbf{Autores}: & Mechoulam, Alan  &  (58438)\\
 & Lambertucci, Guido Enrique  & (58009) \\
 & Rodriguez Turco, Martín Sebastian  & (56629) \\
 & Londero Bonaparte, Tomás Guillermo  & (58150) \\
 &  & \\
 &  & \\
 &  & \\
\textbf{Tutores}: & Orchessi, Walter & \hspace*{4cm} \\
 & Carrillo Del Pino, Manuel Esteban & \hspace*{4cm} \\
 & Gasparini, Ignacio & \hspace*{4cm} \\
 &  & \\
 &  & \\
\textbf{Fecha}: & 06/10/2021 & \hspace*{4.5cm}\\
\end{tabular}

\end{titlepage}

%%%%%%%%%%%%%%%%%%%%%%%%%%%%%%%%%
%			Indice				%
%%%%%%%%%%%%%%%%%%%%%%%%%%%%%%%%%

\MyIndex

\hypersetup{linkcolor=AzulInfo}

\newpage

%%%%%%%%%%%%%%%%%%%%%%%%%%%%%%%%%
%			Informe				%
%%%%%%%%%%%%%%%%%%%%%%%%%%%%%%%%%

\note{Este archivo sirve como guía/tutorial:}

Se crearon varios comandos útiles para desarrollar el informe, todos se encuentran en el archivos \quotes{New-Definitions.tex}. En este mismo archivo se muestran ejemplos de como funcionan la mayoría.

Se sigue una estructura donde cada sección es un archivo distinto. Viendo este documento, el \quotes{Header.tex} y las demás carpetas se puede replicar.

\note{PARA PODER COMPILAR:}
Usar XeLaTex y en la carpeta de Utils instalar el archivo \quotes{AvenirLTStd-Roman.otf}. Actualizar todos los paquetes desde Miktex (si no lo tenés descargalo) y habilitar la instalación automática de los mismos.

\note{Parte del formato:}

USAR \textbackslash Section, \textbackslash Subsection y \textbackslash Subsubsection (existen también las \textbackslash Subsubsubsection) en lugar de \textbackslash section, \textbackslash subsection y \textbackslash subsubsection. De esta forma antes del indice, cada imagen y tabla tendrá el indice de la sección, subsección, subsubsección y subsubsubsección, para identificarlo mejor. Esto se entiende mejor viendo los index de las figuras más abajo.

El paquete de lipsum (linea 13 del header) se puede sacar, es solo para rellenar con texto y probar el formato (borrar también las instancias en este tex).

\observacion{\PrimerObs}{Esto sí se puede ver.}

\observacion{\SegundaObs}{Esto no se puede ver.}

\Section{Acrónimos y definiciones}
\documentclass[a4paper]{article}
%\usepackage[utf8]{inputenc}
\usepackage[spanish, es-tabla, es-noshorthands]{babel}
\usepackage[table,xcdraw,dvipsnames]{xcolor}
\usepackage[a4paper, footnotesep=1.25cm, headheight=1.25cm, top=2.54cm, left=2.54cm,
 bottom=2.54cm, right=2.54cm]{geometry}
%\geometry{showframe}
 \usepackage[normalem]{ulem}
 \useunder{\uline}{\ul}{}

%VERIFICAR EL HEAD Y EL FOOT EN
%https://ctan.dcc.uchile.cl/macros/latex/contrib/geometry/geometry.pdf

%Paquetes varios:
\usepackage{verbatimbox}

%\usepackage{wrapfig}			%Wrap figure in text
\usepackage[export]{adjustbox}	%Move images
\usepackage{changepage}			%Move tables
\usepackage{todonotes}

\usepackage{tikz}
\usepackage{amsmath}
\usepackage{amsfonts}
\usepackage{amssymb}
\usepackage{float}
\usepackage[graphicx]{realboxes}
\usepackage{caption}
\usepackage{subcaption}
\usepackage{multicol}
\usepackage{multirow}
\setlength{\doublerulesep}{\arrayrulewidth}
%\usepackage{booktabs}

\usepackage{array}
\newcolumntype{C}[1]{>{\centering\let\newline\\\arraybackslash\hspace{0pt}}m{#1}}
%\usepackage[american]{circuitikz}
\usetikzlibrary{calc}
\usepackage{fancyhdr}
\usepackage{units} 

\usepackage{colortbl}
%\usepackage{sectsty}
%\usepackage{unicode-math}

%FONTS (IMPORTANTE): Compilar en XeLaTex o LuaLaTeX
\usepackage{anyfontsize}	%Font size
\usepackage{fontspec}		%Font type
%Si sigue sin andar comentar \usepackage[utf8]{inputenc}
%https://ctan.dcc.uchile.cl/macros/unicodetex/latex/fontspec/fontspec.pdf
%https://www.overleaf.com/learn/latex/XeLaTeX

%Path para imagenes para trabajar en subarchivos
\graphicspath{{../Resumen/}{../Referencias/}{../Apendice/}{../Descripción de la Empresa/}{../Tareas del Alumno/}{../Conclusiones/}{../Herramientas Empleadas/}}

%Definiciones de nuevos comandos y colores
%COLORES:
\definecolor{AzulFoot}{rgb}{0.682,0.809,0.926}	%RGB	%{174,206,235}
\definecolor{AzulInfo}{rgb}{0.180,0.455,0.710}	%RGB	%{46,116,181}
\definecolor{AzulTable}{rgb}{0.302,0.507,0.871}	%RGB	%{68,114,196}
\definecolor{PName}{rgb}{0.353,0.353,0.353}		%RGB	%{90,90,90}
\definecolor{mygreen}{rgb}{28,172,0} % color values Red, Green, Blue
\definecolor{mylilas}{rgb}{170,55,241}

%Change Font Size

% #1 = size, #2 = text
\newcommand{\setparagraphsize}[2]{{\fontsize{#1}{6}\selectfont#2 \par}}		%Cambia el size de todo el parrafo
\newcommand{\setlinesize}[2]{{\fontsize{#1}{6}\selectfont#2}}				%Cambia el font de una oración

%IMAGE IN TABLE:			%Ejemplo: \includeintable{.3}{ImagenesFactibilidad/pend}
\renewcommand\fbox{\fcolorbox{white}{white}}
\setlength{\fboxrule}{0pt}	%padding thickness
\setlength{\fboxsep}{4pt}	%border thickness
\newcommand{\includeintable}[2]{	
	\fbox{
		\begin{minipage}{#1\textwidth}
        	\includegraphics[width=\linewidth]{#2}
    	\end{minipage}
	}
}

%LINK IN REF
\newcommand{\reflink}[1]{		%LINK
	\href{#1}{#1}
}

%NOTAS:
\newcommand{\note}[1]{		%RED BIG NOTE (TODO)
	\begin{center}
		\huge{ \textcolor{red}{#1} }
	\end{center}
}

\newcommand{\lnote}[1]{{\fontsize{14}{6}\selectfont\textcolor{green}{#1}}}	%Notas pequeñas

\newcommand{\observacion}[2]{  \ifnumequal{1}{#1}{ { \todo[inline,backgroundcolor=red!25,bordercolor=red!100]{\textbf{Observación: #2}} } }{  }  }

\newcommand{\red}[1]{\textcolor{red}{#1}}

\newcommand{\TBD}{\textcolor{red}{(TBD) }}
\newcommand{\tbd}{\textcolor{red}{(TBD) }}

\newcommand{\TBC}{\textcolor{red}{(TBC) }}
\newcommand{\tbc}{\textcolor{red}{(TBC) }}

\newcommand{\quotes}[1]{``#1''}
\newcommand{\q}[1]{``#1''}

\newcommand{\ip}{192.168.0.10:1880}
\newcommand{\ipadmin}{192.168.0.10:1880/admin}

% Comandos para agregar elementos en tablas de acronimos y definiciones
\newcommand{\addacronym}[2]{\textbf{#1} & \begin{tabular}[l]{@{}l@{}}#2\end{tabular} \\ \hline}

% tabItem
\newcommand{\tabitem}{~~\llap{\textbullet}~~}


\usepackage{hyperref}
\hypersetup{
    colorlinks=true,
    linkcolor=black,
    filecolor=magenta,      
    urlcolor=AzulInfo,
    citecolor=AzulInfo,    
}

%Configuración del header y del footer:
\usepackage{etoolbox}
\pagestyle{fancy}
\fancyhf{}
\rfoot{\thepage}
\renewcommand{\footrulewidth}{4pt}
\renewcommand{\headrulewidth}{0pt}
\patchcmd{\footrule}{\hrule}{\color{AzulFoot}\hrule}{}{}

%Código en el informe
%% IMPORTANTE:
% Verificar que esté \usepackage[dvipsnames]{xcolor}

%\usepackage{listingsutf8}
\usepackage{listings}

\renewcommand{\lstlistingname}{Código}

%LSTSET: Pone un recuadro y contador de linea en el codigo
\newcommand{\boxstyle}{
	\lstset{
		basicstyle=\sffamily\color{black},
		frame=single,
		numbers=left,
		numbersep=5pt,
		numberstyle=\color{gray},
		showspaces=false,
		showstringspaces=false
	}
}

\newcommand{\defaultstyle}{
	\lstset{
		basicstyle=\sffamily\color{white},
		frame=none,
		numbers=none,
		showspaces=true,
		showstringspaces=true
	}
}

\lstdefinelanguage{Kotlin}{
  captionpos=b,
  comment=[l]{//},
  commentstyle={\color{gray}\ttfamily},
  emph={filter, first, firstOrNull, forEach, lazy, map, mapNotNull, println},
  emphstyle={\color{OrangeRed}},
  identifierstyle=\color{black},
  keywords={!in, !is, abstract, actual, annotation, as, as?, break, by, catch, class, companion, const, constructor, continue, crossinline, data, delegate, do, dynamic, else, enum, expect, external, false, field, file, final, finally, for, fun, get, if, import, in, infix, init, inline, inner, interface, internal, is, lateinit, noinline, null, object, open, operator, out, override, package, param, private, property, protected, public, receiveris, reified, return, return@, sealed, set, setparam, super, suspend, tailrec, this, throw, true, try, typealias, typeof, val, var, vararg, when, where, while},
  keywordstyle={\color{NavyBlue}\bfseries},
  morecomment=[s]{/*}{*/},
  morestring=[b]",
  morestring=[s]{"""*}{*"""},
  ndkeywords={@Deprecated, @JvmField, @JvmName, @JvmOverloads, @JvmStatic, @JvmSynthetic, Array, Byte, Double, Float, Int, Integer, Iterable, Long, Runnable, Short, String, Any, Unit, Nothing},
  ndkeywordstyle={\color{BurntOrange}\bfseries},
  sensitive=true,
  stringstyle={\color{ForestGreen}\ttfamily},
}

\lstdefinelanguage{Swift}
{
  morekeywords={
    open,catch,@escaping,nil,throws,func,if,then,else,for,in,while,do,switch,case,default,where,break,continue,fallthrough,return,
    typealias,struct,class,enum,protocol,var,func,let,get,set,willSet,didSet,inout,init,deinit,extension,
    subscript,prefix,operator,infix,postfix,precedence,associativity,left,right,none,convenience,dynamic,
    final,lazy,mutating,nonmutating,optional,override,required,static,unowned,safe,weak,internal,
    private,public,is,as,self,unsafe,dynamicType,true,false,nil,Type,Protocol,
  },
  morecomment=[l]{//}, % l is for line comment
  morecomment=[s]{/*}{*/}, % s is for start and end delimiter
  morestring=[b]", % defines that strings are enclosed in double quotes
  breaklines=true,
  escapeinside={\%*}{*)},
  numbers=left,
  captionpos=b,
  breakatwhitespace=true,
  basicstyle=\linespread{1.0}\ttfamily, % https://tex.stackexchange.com/a/102728/129441
}

\definecolor{keyword}{HTML}{BA2CA3}
\definecolor{string}{HTML}{D12F1B}
\definecolor{comment}{HTML}{008400}

\newcommand{\swiftstyle}{
	\lstset{
  		language=Swift,
  		inputencoding=utf8x,
		extendedchars=\true,
	  	basicstyle=\ttfamily,
	  	showstringspaces=false, % lets spaces in strings appear as real spaces
  		columns=fixed,
  		keepspaces=true,
  		keywordstyle=\color{keyword},
  		stringstyle=\color{string},
  		commentstyle=\color{comment}
	}
}


%Como usarlo:

%\begin{lstlisting}[caption={Simple code listing.}, label={lst:example1}, language=Kotlin]
%// this is a simple code listing:
%println("hello kotlin from latex")
%\end{lstlisting}

%Si se corta en 2 páginas distintas:

%\vspace{1mm}
%\noindent{\begin{minipage}{\linewidth}
%\begin{lstlisting}[...]
%...
%\end{lstlisting}
%\end{minipage}}




\usepackage{titlesec}		%Para hacer las subsubsubsections

%Colores a los nombres de las secciones:
%\sectionfont{\color{AzulInfo}}  % sets color of sections
%\subsectionfont{\color{AzulInfo}}
%\subsubsectionfont{\color{AzulInfo}}

%PICTURES AND TABLE INDEX:
\newcommand{\Section}[1]{ \section{#1} 
	\phantomsection \setcounter{figure}{0} \setcounter{table}{0} \setcounter{lstlisting}{0}
		\renewcommand{\thetable}{\arabic{section}.\arabic{table}}
		\renewcommand{\thefigure}{\arabic{section}.\arabic{figure}}
		\renewcommand{\thelstlisting}{\arabic{section}.\arabic{lstlisting}}
}

\newcommand{\Subsection}[1]{ \subsection{#1}
	\phantomsection \setcounter{figure}{0} \setcounter{table}{0} \setcounter{lstlisting}{0}
		\renewcommand{\thetable}{\arabic{section}.\arabic{subsection}.\arabic{table}}
		\renewcommand{\thefigure}{\arabic{section}.\arabic{subsection}.\arabic{figure}}
		\renewcommand{\thelstlisting}{\arabic{section}.\arabic{subsection}.\arabic{lstlisting}}
}

\newcommand{\Subsubsection}[1]{ \subsubsection{#1} 
	\phantomsection \setcounter{figure}{0} \setcounter{table}{0}  \setcounter{lstlisting}{0}
		\renewcommand{\thetable}{\arabic{section}.\arabic{subsection}.\arabic{subsubsection}.\arabic{table}}
		\renewcommand{\thefigure}{\arabic{section}.\arabic{subsection}.\arabic{subsubsection}.\arabic{figure}}
		\renewcommand{\thelstlisting}{\arabic{section}.\arabic{subsection}.\arabic{subsubsection}.\arabic{lstlisting}}
}

%Definición de subsubsubsection:
\titleclass{\subsubsubsection}{straight}[\subsection]

\newcounter{subsubsubsection}[subsubsection]
\renewcommand\thesubsubsubsection{\thesubsubsection.\arabic{subsubsubsection}}

\titleformat{\subsubsubsection}
  {\normalfont\normalsize\bfseries\color{AzulInfo}}{\thesubsubsubsection}{1em}{}	%Color de subsubsubsection
\titlespacing*{\subsubsubsection}
{0pt}{3.25ex plus 1ex minus .2ex}{1.5ex plus .2ex}

\makeatletter
\renewcommand\paragraph{\@startsection{paragraph}{5}{\z@}%
  {3.25ex \@plus1ex \@minus.2ex}%
  {-1em}%
  {\normalfont\normalsize\bfseries}}
\renewcommand\subparagraph{\@startsection{subparagraph}{6}{\parindent}%
  {3.25ex \@plus1ex \@minus .2ex}%
  {-1em}%
  {\normalfont\normalsize\bfseries}}
\def\toclevel@subsubsubsection{4}
\def\toclevel@paragraph{5}
\def\toclevel@paragraph{6}
\def\l@subsubsubsection{\@dottedtocline{4}{7em}{4em}}
\def\l@paragraph{\@dottedtocline{5}{10em}{5em}}
\def\l@subparagraph{\@dottedtocline{6}{14em}{6em}}
\makeatother

\setcounter{secnumdepth}{4}
\setcounter{tocdepth}{4}

%Subsubsubsection:
\newcommand{\Subsubsubsection}[1]{ \subsubsubsection{#1} 
	\phantomsection \setcounter{figure}{0} \setcounter{table}{0} \renewcommand{\thetable}{\arabic{section}.\arabic{subsection}.\arabic{subsubsection}.\arabic{subsubsubsection}.\arabic{table}} \renewcommand{\thefigure}{\arabic{section}.\arabic{subsection}.\arabic{subsubsection}.\arabic{subsubsubsection}.\arabic{figure}}
}

%Tamaño, color e identación de sección, subsección, subsubsección y subsubsubsección:
%La identación de las subsecciones está tambien en Index-cfg.tex para el toc, lot y lot en el index
\titleformat{\section}[block]{\fontsize{16}{6}\selectfont\bfseries\color{AzulInfo}}{\thesection.}{1em}{} 
\titleformat{\subsection}[block]{\hspace{2.5em}\fontsize{13}{6}\selectfont\color{AzulInfo}}{\thesubsection}{1em}{}
\titleformat{\subsubsection}[block]{\hspace{3.5em}\fontsize{12}{6}\selectfont\color{AzulInfo}}{\thesubsubsection}{1em}{}
\titleformat{\subsubsubsection}[block]{\hspace{4em}\fontsize{11}{6}\selectfont\color{AzulInfo}}{\thesubsubsubsection}{1em}{}

%Pone las refrencias en el indice
\usepackage[numbib, nottoc, notlot, notlof]{tocbibind}

%Pone toc, lof y lot en colores y elijo el titulo de estos
\addto\captionsspanish{
	\renewcommand\contentsname{Contenidos}
	\renewcommand\listfigurename{Lista de Figuras}
	\renewcommand\listtablename{Lista de Tablas}
}

%Agrega TOC al indice
\renewcommand{\tableofcontents}{
	\stepcounter{section}
	\addcontentsline{toc}{section}{\protect\numberline{\thesection}\textbf{Contenidos}}
	\tableofcontents
}

%Agrega LOF al indice
\renewcommand{\listoffigures}{
	\stepcounter{section}
	\addcontentsline{toc}{section}{\protect\numberline{\thesection}\textbf{Lista de Figuras}}
	\listoffigures
}

%Agrega LOT al indice
\renewcommand{\listoftables}{
	\stepcounter{section}
	\addcontentsline{toc}{section}{\protect\numberline{\thesection}\textbf{Lista de Tablas}}
	\listoftables
}

%Indices: cambio la separación de los numeros para que entren tablas y fotos
\usepackage{tocloft}
\setlength{\cftfignumwidth}{1.35cm}  % change numwidth from figures in lof
\setlength{\cfttabnumwidth}{1.35cm}  % change numwidth from tables in lot
\renewcommand{\cfttoctitlefont}{\Large\bfseries\color{AzulInfo}}
\renewcommand{\cftloftitlefont}{\Large\bfseries\color{AzulInfo}}
\renewcommand{\cftlottitlefont}{\Large\bfseries\color{AzulInfo}}

%Coloca lineas punteadas a las seciones en el TOC
\renewcommand{\cftsecleader}{\cftdotfill{\cftdotsep}}

%Items con bullets y no cuadrados
\renewcommand{\labelitemi}{\textbullet }


\begin{document}

\begin{table}[H]
\centering
\begin{tabular}{!{\color{AzulTable}\vrule}ll!{\color{AzulTable}\vrule}}
\arrayrulecolor{AzulTable}
\hline
\rowcolor{AzulTable}
\multicolumn{1}{|c}{\textcolor{white}{Acrónimo}} & \multicolumn{1}{c|}{\textcolor{white}{Descripción}} \\ \hline
\textbf{AMB}                         & Ambiente                         \\
\textbf{COM}                         & Comunicación                         \\
\textbf{EGM}                         & Electromecánica                         \\
\textbf{MEC}                         & Mecánico                         \\
\textbf{SER}                         & Servicio Técnico                         \\
\textbf{TC}                         & Termocupla                         \\
\textbf{IC}                         & Circuito Integrado                         \\
\textbf{RTD}                         & Resistance Temperature Detector           \\	\hline
\end{tabular}
\end{table}

\begin{table}[H]
\centering
\begin{tabular}{!{\color{AzulTable}\vrule}ll!{\color{AzulTable}\vrule}}
\arrayrulecolor{AzulTable}
\hline
\rowcolor{AzulTable}
\multicolumn{1}{|c}{\textcolor{white}{Término}} & \multicolumn{1}{c|}{\textcolor{white}{Descripción}} \\ \hline
\textbf{1}                         & 2                         \\
\textbf{3}                         & 4                         \\ \hline
\end{tabular}
\end{table}

\end{document}

\Section{Lorem Ipsum}

\lipsum[1-2]
																			
\setparagraphsize{16}{\textcolor{red}{\lipsum[3-4]}}

También se puede \setlinesize{6}{cambiar el tamaño} en el medio de una oración. Pequeñas \lnote{notas verdes} para poner una nota (y que se vea) en el texto \cite{ref:cita1}.

Para la versión del informe que se entregue se puede usar \tbd y \tbc: El producto costará \TBD, ocupará un volumen de \tbd, pesará \TBC y durará \tbc.

\note{Nota más importante para no olvidar.}

\Section{Ejercicico N}
\documentclass[a4paper]{article}
%\usepackage[utf8]{inputenc}
\usepackage[spanish, es-tabla, es-noshorthands]{babel}
\usepackage[table,xcdraw,dvipsnames]{xcolor}
\usepackage[a4paper, footnotesep=1.25cm, headheight=1.25cm, top=2.54cm, left=2.54cm,
 bottom=2.54cm, right=2.54cm]{geometry}
%\geometry{showframe}
 \usepackage[normalem]{ulem}
 \useunder{\uline}{\ul}{}

%VERIFICAR EL HEAD Y EL FOOT EN
%https://ctan.dcc.uchile.cl/macros/latex/contrib/geometry/geometry.pdf

%Paquetes varios:
\usepackage{verbatimbox}

%\usepackage{wrapfig}			%Wrap figure in text
\usepackage[export]{adjustbox}	%Move images
\usepackage{changepage}			%Move tables
\usepackage{todonotes}

\usepackage{tikz}
\usepackage{amsmath}
\usepackage{amsfonts}
\usepackage{amssymb}
\usepackage{float}
\usepackage[graphicx]{realboxes}
\usepackage{caption}
\usepackage{subcaption}
\usepackage{multicol}
\usepackage{multirow}
\setlength{\doublerulesep}{\arrayrulewidth}
%\usepackage{booktabs}

\usepackage{array}
\newcolumntype{C}[1]{>{\centering\let\newline\\\arraybackslash\hspace{0pt}}m{#1}}
%\usepackage[american]{circuitikz}
\usetikzlibrary{calc}
\usepackage{fancyhdr}
\usepackage{units} 

\usepackage{colortbl}
%\usepackage{sectsty}
%\usepackage{unicode-math}

%FONTS (IMPORTANTE): Compilar en XeLaTex o LuaLaTeX
\usepackage{anyfontsize}	%Font size
\usepackage{fontspec}		%Font type
%Si sigue sin andar comentar \usepackage[utf8]{inputenc}
%https://ctan.dcc.uchile.cl/macros/unicodetex/latex/fontspec/fontspec.pdf
%https://www.overleaf.com/learn/latex/XeLaTeX

%Path para imagenes para trabajar en subarchivos
\graphicspath{{../Resumen/}{../Referencias/}{../Apendice/}{../Descripción de la Empresa/}{../Tareas del Alumno/}{../Conclusiones/}{../Herramientas Empleadas/}}

%Definiciones de nuevos comandos y colores
%COLORES:
\definecolor{AzulFoot}{rgb}{0.682,0.809,0.926}	%RGB	%{174,206,235}
\definecolor{AzulInfo}{rgb}{0.180,0.455,0.710}	%RGB	%{46,116,181}
\definecolor{AzulTable}{rgb}{0.302,0.507,0.871}	%RGB	%{68,114,196}
\definecolor{PName}{rgb}{0.353,0.353,0.353}		%RGB	%{90,90,90}
\definecolor{mygreen}{rgb}{28,172,0} % color values Red, Green, Blue
\definecolor{mylilas}{rgb}{170,55,241}

%Change Font Size

% #1 = size, #2 = text
\newcommand{\setparagraphsize}[2]{{\fontsize{#1}{6}\selectfont#2 \par}}		%Cambia el size de todo el parrafo
\newcommand{\setlinesize}[2]{{\fontsize{#1}{6}\selectfont#2}}				%Cambia el font de una oración

%IMAGE IN TABLE:			%Ejemplo: \includeintable{.3}{ImagenesFactibilidad/pend}
\renewcommand\fbox{\fcolorbox{white}{white}}
\setlength{\fboxrule}{0pt}	%padding thickness
\setlength{\fboxsep}{4pt}	%border thickness
\newcommand{\includeintable}[2]{	
	\fbox{
		\begin{minipage}{#1\textwidth}
        	\includegraphics[width=\linewidth]{#2}
    	\end{minipage}
	}
}

%LINK IN REF
\newcommand{\reflink}[1]{		%LINK
	\href{#1}{#1}
}

%NOTAS:
\newcommand{\note}[1]{		%RED BIG NOTE (TODO)
	\begin{center}
		\huge{ \textcolor{red}{#1} }
	\end{center}
}

\newcommand{\lnote}[1]{{\fontsize{14}{6}\selectfont\textcolor{green}{#1}}}	%Notas pequeñas

\newcommand{\observacion}[2]{  \ifnumequal{1}{#1}{ { \todo[inline,backgroundcolor=red!25,bordercolor=red!100]{\textbf{Observación: #2}} } }{  }  }

\newcommand{\red}[1]{\textcolor{red}{#1}}

\newcommand{\TBD}{\textcolor{red}{(TBD) }}
\newcommand{\tbd}{\textcolor{red}{(TBD) }}

\newcommand{\TBC}{\textcolor{red}{(TBC) }}
\newcommand{\tbc}{\textcolor{red}{(TBC) }}

\newcommand{\quotes}[1]{``#1''}
\newcommand{\q}[1]{``#1''}

\newcommand{\ip}{192.168.0.10:1880}
\newcommand{\ipadmin}{192.168.0.10:1880/admin}

% Comandos para agregar elementos en tablas de acronimos y definiciones
\newcommand{\addacronym}[2]{\textbf{#1} & \begin{tabular}[l]{@{}l@{}}#2\end{tabular} \\ \hline}

% tabItem
\newcommand{\tabitem}{~~\llap{\textbullet}~~}


\usepackage{hyperref}
\hypersetup{
    colorlinks=true,
    linkcolor=black,
    filecolor=magenta,      
    urlcolor=AzulInfo,
    citecolor=AzulInfo,    
}

%Configuración del header y del footer:
\usepackage{etoolbox}
\pagestyle{fancy}
\fancyhf{}
\rfoot{\thepage}
\renewcommand{\footrulewidth}{4pt}
\renewcommand{\headrulewidth}{0pt}
\patchcmd{\footrule}{\hrule}{\color{AzulFoot}\hrule}{}{}

%Código en el informe
%% IMPORTANTE:
% Verificar que esté \usepackage[dvipsnames]{xcolor}

%\usepackage{listingsutf8}
\usepackage{listings}

\renewcommand{\lstlistingname}{Código}

%LSTSET: Pone un recuadro y contador de linea en el codigo
\newcommand{\boxstyle}{
	\lstset{
		basicstyle=\sffamily\color{black},
		frame=single,
		numbers=left,
		numbersep=5pt,
		numberstyle=\color{gray},
		showspaces=false,
		showstringspaces=false
	}
}

\newcommand{\defaultstyle}{
	\lstset{
		basicstyle=\sffamily\color{white},
		frame=none,
		numbers=none,
		showspaces=true,
		showstringspaces=true
	}
}

\lstdefinelanguage{Kotlin}{
  captionpos=b,
  comment=[l]{//},
  commentstyle={\color{gray}\ttfamily},
  emph={filter, first, firstOrNull, forEach, lazy, map, mapNotNull, println},
  emphstyle={\color{OrangeRed}},
  identifierstyle=\color{black},
  keywords={!in, !is, abstract, actual, annotation, as, as?, break, by, catch, class, companion, const, constructor, continue, crossinline, data, delegate, do, dynamic, else, enum, expect, external, false, field, file, final, finally, for, fun, get, if, import, in, infix, init, inline, inner, interface, internal, is, lateinit, noinline, null, object, open, operator, out, override, package, param, private, property, protected, public, receiveris, reified, return, return@, sealed, set, setparam, super, suspend, tailrec, this, throw, true, try, typealias, typeof, val, var, vararg, when, where, while},
  keywordstyle={\color{NavyBlue}\bfseries},
  morecomment=[s]{/*}{*/},
  morestring=[b]",
  morestring=[s]{"""*}{*"""},
  ndkeywords={@Deprecated, @JvmField, @JvmName, @JvmOverloads, @JvmStatic, @JvmSynthetic, Array, Byte, Double, Float, Int, Integer, Iterable, Long, Runnable, Short, String, Any, Unit, Nothing},
  ndkeywordstyle={\color{BurntOrange}\bfseries},
  sensitive=true,
  stringstyle={\color{ForestGreen}\ttfamily},
}

\lstdefinelanguage{Swift}
{
  morekeywords={
    open,catch,@escaping,nil,throws,func,if,then,else,for,in,while,do,switch,case,default,where,break,continue,fallthrough,return,
    typealias,struct,class,enum,protocol,var,func,let,get,set,willSet,didSet,inout,init,deinit,extension,
    subscript,prefix,operator,infix,postfix,precedence,associativity,left,right,none,convenience,dynamic,
    final,lazy,mutating,nonmutating,optional,override,required,static,unowned,safe,weak,internal,
    private,public,is,as,self,unsafe,dynamicType,true,false,nil,Type,Protocol,
  },
  morecomment=[l]{//}, % l is for line comment
  morecomment=[s]{/*}{*/}, % s is for start and end delimiter
  morestring=[b]", % defines that strings are enclosed in double quotes
  breaklines=true,
  escapeinside={\%*}{*)},
  numbers=left,
  captionpos=b,
  breakatwhitespace=true,
  basicstyle=\linespread{1.0}\ttfamily, % https://tex.stackexchange.com/a/102728/129441
}

\definecolor{keyword}{HTML}{BA2CA3}
\definecolor{string}{HTML}{D12F1B}
\definecolor{comment}{HTML}{008400}

\newcommand{\swiftstyle}{
	\lstset{
  		language=Swift,
  		inputencoding=utf8x,
		extendedchars=\true,
	  	basicstyle=\ttfamily,
	  	showstringspaces=false, % lets spaces in strings appear as real spaces
  		columns=fixed,
  		keepspaces=true,
  		keywordstyle=\color{keyword},
  		stringstyle=\color{string},
  		commentstyle=\color{comment}
	}
}


%Como usarlo:

%\begin{lstlisting}[caption={Simple code listing.}, label={lst:example1}, language=Kotlin]
%// this is a simple code listing:
%println("hello kotlin from latex")
%\end{lstlisting}

%Si se corta en 2 páginas distintas:

%\vspace{1mm}
%\noindent{\begin{minipage}{\linewidth}
%\begin{lstlisting}[...]
%...
%\end{lstlisting}
%\end{minipage}}




\usepackage{titlesec}		%Para hacer las subsubsubsections

%Colores a los nombres de las secciones:
%\sectionfont{\color{AzulInfo}}  % sets color of sections
%\subsectionfont{\color{AzulInfo}}
%\subsubsectionfont{\color{AzulInfo}}

%PICTURES AND TABLE INDEX:
\newcommand{\Section}[1]{ \section{#1} 
	\phantomsection \setcounter{figure}{0} \setcounter{table}{0} \setcounter{lstlisting}{0}
		\renewcommand{\thetable}{\arabic{section}.\arabic{table}}
		\renewcommand{\thefigure}{\arabic{section}.\arabic{figure}}
		\renewcommand{\thelstlisting}{\arabic{section}.\arabic{lstlisting}}
}

\newcommand{\Subsection}[1]{ \subsection{#1}
	\phantomsection \setcounter{figure}{0} \setcounter{table}{0} \setcounter{lstlisting}{0}
		\renewcommand{\thetable}{\arabic{section}.\arabic{subsection}.\arabic{table}}
		\renewcommand{\thefigure}{\arabic{section}.\arabic{subsection}.\arabic{figure}}
		\renewcommand{\thelstlisting}{\arabic{section}.\arabic{subsection}.\arabic{lstlisting}}
}

\newcommand{\Subsubsection}[1]{ \subsubsection{#1} 
	\phantomsection \setcounter{figure}{0} \setcounter{table}{0}  \setcounter{lstlisting}{0}
		\renewcommand{\thetable}{\arabic{section}.\arabic{subsection}.\arabic{subsubsection}.\arabic{table}}
		\renewcommand{\thefigure}{\arabic{section}.\arabic{subsection}.\arabic{subsubsection}.\arabic{figure}}
		\renewcommand{\thelstlisting}{\arabic{section}.\arabic{subsection}.\arabic{subsubsection}.\arabic{lstlisting}}
}

%Definición de subsubsubsection:
\titleclass{\subsubsubsection}{straight}[\subsection]

\newcounter{subsubsubsection}[subsubsection]
\renewcommand\thesubsubsubsection{\thesubsubsection.\arabic{subsubsubsection}}

\titleformat{\subsubsubsection}
  {\normalfont\normalsize\bfseries\color{AzulInfo}}{\thesubsubsubsection}{1em}{}	%Color de subsubsubsection
\titlespacing*{\subsubsubsection}
{0pt}{3.25ex plus 1ex minus .2ex}{1.5ex plus .2ex}

\makeatletter
\renewcommand\paragraph{\@startsection{paragraph}{5}{\z@}%
  {3.25ex \@plus1ex \@minus.2ex}%
  {-1em}%
  {\normalfont\normalsize\bfseries}}
\renewcommand\subparagraph{\@startsection{subparagraph}{6}{\parindent}%
  {3.25ex \@plus1ex \@minus .2ex}%
  {-1em}%
  {\normalfont\normalsize\bfseries}}
\def\toclevel@subsubsubsection{4}
\def\toclevel@paragraph{5}
\def\toclevel@paragraph{6}
\def\l@subsubsubsection{\@dottedtocline{4}{7em}{4em}}
\def\l@paragraph{\@dottedtocline{5}{10em}{5em}}
\def\l@subparagraph{\@dottedtocline{6}{14em}{6em}}
\makeatother

\setcounter{secnumdepth}{4}
\setcounter{tocdepth}{4}

%Subsubsubsection:
\newcommand{\Subsubsubsection}[1]{ \subsubsubsection{#1} 
	\phantomsection \setcounter{figure}{0} \setcounter{table}{0} \renewcommand{\thetable}{\arabic{section}.\arabic{subsection}.\arabic{subsubsection}.\arabic{subsubsubsection}.\arabic{table}} \renewcommand{\thefigure}{\arabic{section}.\arabic{subsection}.\arabic{subsubsection}.\arabic{subsubsubsection}.\arabic{figure}}
}

%Tamaño, color e identación de sección, subsección, subsubsección y subsubsubsección:
%La identación de las subsecciones está tambien en Index-cfg.tex para el toc, lot y lot en el index
\titleformat{\section}[block]{\fontsize{16}{6}\selectfont\bfseries\color{AzulInfo}}{\thesection.}{1em}{} 
\titleformat{\subsection}[block]{\hspace{2.5em}\fontsize{13}{6}\selectfont\color{AzulInfo}}{\thesubsection}{1em}{}
\titleformat{\subsubsection}[block]{\hspace{3.5em}\fontsize{12}{6}\selectfont\color{AzulInfo}}{\thesubsubsection}{1em}{}
\titleformat{\subsubsubsection}[block]{\hspace{4em}\fontsize{11}{6}\selectfont\color{AzulInfo}}{\thesubsubsubsection}{1em}{}

%Pone las refrencias en el indice
\usepackage[numbib, nottoc, notlot, notlof]{tocbibind}

%Pone toc, lof y lot en colores y elijo el titulo de estos
\addto\captionsspanish{
	\renewcommand\contentsname{Contenidos}
	\renewcommand\listfigurename{Lista de Figuras}
	\renewcommand\listtablename{Lista de Tablas}
}

%Agrega TOC al indice
\renewcommand{\tableofcontents}{
	\stepcounter{section}
	\addcontentsline{toc}{section}{\protect\numberline{\thesection}\textbf{Contenidos}}
	\tableofcontents
}

%Agrega LOF al indice
\renewcommand{\listoffigures}{
	\stepcounter{section}
	\addcontentsline{toc}{section}{\protect\numberline{\thesection}\textbf{Lista de Figuras}}
	\listoffigures
}

%Agrega LOT al indice
\renewcommand{\listoftables}{
	\stepcounter{section}
	\addcontentsline{toc}{section}{\protect\numberline{\thesection}\textbf{Lista de Tablas}}
	\listoftables
}

%Indices: cambio la separación de los numeros para que entren tablas y fotos
\usepackage{tocloft}
\setlength{\cftfignumwidth}{1.35cm}  % change numwidth from figures in lof
\setlength{\cfttabnumwidth}{1.35cm}  % change numwidth from tables in lot
\renewcommand{\cfttoctitlefont}{\Large\bfseries\color{AzulInfo}}
\renewcommand{\cftloftitlefont}{\Large\bfseries\color{AzulInfo}}
\renewcommand{\cftlottitlefont}{\Large\bfseries\color{AzulInfo}}

%Coloca lineas punteadas a las seciones en el TOC
\renewcommand{\cftsecleader}{\cftdotfill{\cftdotsep}}

%Items con bullets y no cuadrados
\renewcommand{\labelitemi}{\textbullet }


\begin{document}

\TBD

\end{document}

\newpage
\begin{flushleft}
\begin{thebibliography}{9}
\bibitem{ref:cita1}
T. Bangalter, G. M. de Homem-Christo. \textit{Random Acces Memory}. Vol. 8, 2013.

\end{thebibliography}
\end{flushleft}


\end{document}
